% Informe de contexto regional de DespachoX (equipo 19, Hackathon UNAL Manizales 2026).
% Documento vivo. Cada sección vive en secciones/ y se une aquí con \input.
% Compilar con ./compilar.sh (lualatex mediante latexmk).

\documentclass[11pt,a4paper]{article}

% ---------- Datos de la versión (cambiar aquí al publicar una versión nueva)
\newcommand{\versioninforme}{0.1}
\newcommand{\fechaversion}{viernes 20 de noviembre de 2026}

% ---------- Idioma y fuentes
\usepackage{fontspec}
\usepackage{babel}
\babelprovide[import, main]{spanish}
\setmainfont{TeX Gyre Pagella}
\setsansfont{TeX Gyre Heros}
\setmonofont{Latin Modern Mono}[Scale=0.95]

% ---------- Página y tipografía
\usepackage[a4paper, top=2.6cm, bottom=2.6cm, left=2.5cm, right=2.5cm, headheight=14pt]{geometry}
\usepackage[protrusion=true, expansion=false]{microtype}
% No hay patrones de guionado para el español instalados: texto en bandera
% (ragged2e) y sin guiones, para evitar espacios feos o cortes con reglas de otro idioma.
\usepackage[document]{ragged2e}
\hyphenpenalty=10000
\exhyphenpenalty=10000
\emergencystretch=2em
\setlength{\parskip}{0.45em}
\setlength{\RaggedRightParindent}{0pt}

\usepackage{xcolor}
\definecolor{titulo}{RGB}{31,58,84}
\definecolor{marca}{RGB}{150,75,0}
\definecolor{gris}{RGB}{95,95,95}

\usepackage{titlesec}
\titleformat{\section}{\sffamily\Large\bfseries\color{titulo}}{\thesection.}{0.5em}{}
\titleformat{\subsection}{\sffamily\large\bfseries\color{titulo}}{\thesubsection}{0.6em}{}
\titleformat{\paragraph}[runin]{\sffamily\bfseries\color{titulo}}{}{0pt}{}[.]
\titlespacing*{\section}{0pt}{1.6em}{0.7em}
\titlespacing*{\subsection}{0pt}{1.2em}{0.4em}
\titlespacing*{\paragraph}{0pt}{0.7em}{0.6em}

\usepackage{enumitem}
\setlist{itemsep=0.2em, topsep=0.3em, leftmargin=1.4em}

\usepackage{fancyhdr}
\pagestyle{fancy}
\fancyhf{}
\fancyhead[L]{\sffamily\footnotesize\color{gris} DespachoX · Informe de contexto regional}
\fancyhead[R]{\sffamily\footnotesize\color{gris} Versión \versioninforme}
\fancyfoot[C]{\sffamily\footnotesize\color{gris}\thepage}
\renewcommand{\headrulewidth}{0.3pt}

% ---------- Tablas
\usepackage{array}
\usepackage{longtable}
\usepackage{booktabs}
\newcolumntype{L}[1]{>{\RaggedRight\arraybackslash}p{#1}}
\setlength{\LTpre}{0.6em}
\setlength{\LTpost}{0.6em}
\usepackage{caption}
\captionsetup{font={small,sf}, labelfont=bf, labelsep=period, justification=raggedright, singlelinecheck=false}

% ---------- Enlaces
\usepackage{xurl}
\usepackage{hyperref}
\hypersetup{
  colorlinks=true,
  linkcolor=titulo,
  urlcolor=titulo,
  pdftitle={Manizales después del sismo: llevar la comida a cada barrio y vereda},
  pdfauthor={Equipo 19, DespachoX},
  pdfsubject={Informe de contexto regional, versión \versioninforme},
  pdflang={es-CO}
}
\urlstyle{same}
% Rutas de archivos del repositorio, con cortes de línea permitidos.
\DeclareUrlCommand\ruta{\urlstyle{tt}\def\UrlBreaks{\do\/\do\_}\def\UrlBigBreaks{}}

% ---------- Marcas del informe
\newcommand{\pendiente}{\textsf{\small\bfseries\color{marca}pendiente}}
\newcommand{\eninvestigacion}{\textsf{\small\bfseries\color{marca}en investigación}}
\newcommand{\porverificar}{\textsf{\small\bfseries\color{marca}por verificar}}

% ---------- Fecha y hora de compilación, en hora de Colombia (UTC-5, sin horario de verano)
\newcommand{\horacompilacion}{\directlua{tex.sprint(os.date(string.char(33,37).."H:"..string.char(37).."M", os.time() - 5*3600))}}
\newcommand{\fechacompilacion}{\directlua{tex.sprint(os.date(string.char(33,37).."d/"..string.char(37).."m/"..string.char(37).."Y", os.time() - 5*3600))}}

\setcounter{tocdepth}{1}

\begin{document}
\renewcommand{\contentsname}{Contenido}
\renewcommand{\tablename}{Tabla}

% Portada
\begin{titlepage}
\RaggedRight
\vspace*{2.2cm}
{\sffamily\small\color{gris} HACKATHON UNAL MANIZALES 2026\quad·\quad VERTICAL 04, LOGÍSTICA Y ABASTECIMIENTO\par}
\vspace{1.4cm}
{\sffamily\bfseries\fontsize{25}{31}\selectfont\color{titulo} Manizales después del sismo: llevar la comida a cada barrio y vereda\par}
\vspace{0.9cm}
{\Large Informe de contexto regional de DespachoX. Documento vivo\par}
\vspace{2.4cm}
{\sffamily\large Equipo 19, Hackathon UNAL Manizales 2026\par}
\vspace{0.5cm}
{\large Versión \versioninforme, \fechaversion\par}
\vspace{0.2cm}
{\color{gris} Compilado el \fechacompilacion{} a las \horacompilacion, hora de Colombia\par}
\vfill
{\color{titulo}\rule{\linewidth}{0.5pt}}\par
\vspace{0.3em}
{\small Estado de las cifras. Ninguna cifra de este informe ha sido verificada todavía por una persona del equipo. La sección 1 explica los niveles de evidencia y la regla para llevar una cifra al pitch.\par}
\end{titlepage}
\setcounter{page}{2}

% Resumen (150 a 250 palabras, tres o cuatro cifras como mucho)
\section*{Resumen}
\addcontentsline{toc}{section}{Resumen}

Este informe reúne el contexto regional del problema que trabaja DespachoX, pensado para quien no vive en Manizales. El 10 de agosto de 2026, un sismo de magnitud 7,4 con epicentro en el Chocó llevó al Gobierno a declarar la emergencia en Caldas y en otros departamentos. En Manizales, el censo que cerró el 24 de agosto registró 8.240 personas damnificadas.

La comida llegó a la ciudad. La red de bancos de alimentos usó a Manizales como punto de acopio y, sumando todo su corredor, movió más de 2.000 toneladas en la primera semana. La dificultad apareció al repartir. Dos días después seguían cerradas la vía Pereira–Chinchiná y la vía a Fresno por el Alto de Letras. Los cortes de energía dañaron la cadena de frío, y la red señaló la llegada a las zonas rurales como su mayor reto.

DespachoX propone una mesa de despacho que recibe por chat los pedidos de los barrios, recalcula rutas cuando un conductor reporta un cierre y sigue cada paquete hasta su entrega. Una persona aprueba cada plan. En simulación, el sistema supera con claridad al despacho por orden de llegada en entregas a tiempo a puntos críticos. Frente a un despachador con oficio empata en ese indicador, así que el equipo sitúa su aporte en poner esa pericia al alcance de cualquier voluntario.

Ninguna cifra ha sido verificada todavía por una persona del equipo. La pregunta abierta que más pesa es cómo despacha hoy el Banco de Alimentos de Manizales.

\vspace{1.2em}
{\small\tableofcontents}
\clearpage

% Sección 1
\section{Para qué este informe y cómo leerlo}

Este informe responde a un encargo de la mentoría de la hackathon. Los mentores nos pidieron conocer el problema a fondo antes de defender una solución, y dejaron cinco preguntas. Quieren saber por qué hacemos esto y para qué sirve, y a quién le sirve y a quién impacta. También preguntaron si existe algo parecido y por qué no se usa en Manizales o en la región. Las dos últimas miran hacia adelante y piden el prototipo que medirá el efecto en la región y el piloto de 30 a 90 días. La expresión que pusieron en el centro fue «recuperación del territorio». Con ella señalaron que la respuesta tiene que hablar de personas y de lugares además de tecnología.

El documento está escrito para quien no vive en la región. Reúne lo que el equipo ya investigó y lo ordena alrededor de esas preguntas. Es un documento vivo. Esta versión 0.1 recoge los informes de investigación que hay en el repositorio, y las siguientes añadirán lo que traigan otros agentes durante la noche.

\subsection{Niveles de evidencia}

Cada cifra y cada afirmación del informe lleva uno de cuatro niveles de evidencia. Las tablas los abrevian con la palabra en negrilla.

\begin{description}[leftmargin=1.4em, style=nextline, itemsep=0.35em]
\item[\textsf{Oficial}.] Fuente oficial. Es un documento de la entidad que produce el dato o un decreto. Por ejemplo, el Decreto 1261 de 2026 cita el balance de la Unidad Nacional para la Gestión del Riesgo (UNGRD).
\item[\textsf{Prensa}.] Prensa que cita a una entidad. Es una nota de prensa, o el documento de un tercero, que reproduce lo que dijo una entidad o un funcionario. Una cifra que la investigación del equipo calificó con confianza media queda en este nivel aunque venga de un documento institucional, como los informes de Asocapitales. Así no subimos el nivel que le dio la investigación.
\item[\textsf{Inferencia}.] Inferencia nuestra. Es una conclusión del equipo a partir de los datos, o un resultado de nuestra simulación. Las cifras de simulación salen de un modelo con supuestos propios y no de una operación real.
\item[\textsf{Hipótesis}.] Hipótesis por validar. Es algo que creemos y que nadie ha confirmado todavía. Se valida hablando con quien estuvo en la operación.
\end{description}

Las tablas añaden entre paréntesis la confianza que asignó el informe de investigación del equipo, sin cambios. «Alta» corresponde a una entidad oficial o un decreto. «Media» corresponde a prensa con fuente citada o a una cifra preliminar. Cuando el dato viene de un informe que no calificó la confianza, la tabla dice «sin calificar».

\subsection{La regla para llevar una cifra al pitch}

El equipo tiene una regla. Ninguna cifra va al pitch hasta que una persona abra la fuente y la anote en \ruta{docs/FUENTES.md}. Hoy ninguna cifra de este informe cumple esa condición. Los datos vienen de informes que prepararon agentes de investigación el 20 de noviembre de 2026, y nadie del equipo ha abierto todavía los enlaces. Varias consultas web de esos agentes quedaron bloqueadas, de modo que hay vacíos que este informe señala sin rellenarlos.

Las cifras del sismo son preliminares y cambian con cada corte. Por eso cada una aparece con su fecha, y cuando dos fuentes difieren se dan las dos. El informe usa además tres marcas. \pendiente{} indica una respuesta que nadie del equipo ha conseguido todavía. \eninvestigacion{} indica algo que dos agentes están buscando esta noche. \porverificar{} señala un dato que nos pareció dudoso y que dejamos tal como está en la fuente.

Las fuentes de cada dato están en los informes de investigación del repositorio, en la carpeta \ruta{docs/investigacion/}. Los que más usa este informe son el de evidencia del problema (\ruta{04_evidencia_del_problema.md}) y el de referentes latinoamericanos (\ruta{02_referentes_latinoamerica.md}). El anexo reúne en una tabla las fuentes que usa este informe.

% Sección 2
\section{Qué pasó}

\subsection{El sismo y la emergencia}

El sismo ocurrió el 10 de agosto de 2026 a las 7:34 de la mañana. Tuvo magnitud 7,4 y epicentro en San José del Palmar, en el Chocó, según el boletín del Servicio Geológico Colombiano que citan los considerandos del Decreto 1261. El mismo decreto da dos profundidades distintas, 103,4 km en un considerando y unos 94 km en otro. El 19 de agosto, el Gobierno declaró el estado de emergencia por 30 días calendario en quince departamentos, entre ellos Caldas.

El balance nacional creció con los días. El decreto cita el reporte 36 de la UNGRD, con corte al 19 de agosto. Ese reporte contaba 314 personas fallecidas, 4.437 heridas y 262 desaparecidas, y 262.154 damnificados en 147.464 familias. Seis días después, la Organización Panamericana de la Salud (OPS) citaba otro corte de la UNGRD, con 331 fallecidos y 361.707 afectados. Las cifras tempranas eran más bajas. El 10 de agosto Asocapitales contaba 87 muertos en las capitales, y el 12 de agosto la UNGRD reportaba 241, según la Federación Internacional de la Cruz Roja (IFRC). En octubre, la Defensoría hablaba de 335 muertos en 500 municipios de 16 departamentos.

La cobertura territorial también cambió de un corte a otro. El decreto habla de 471 municipios con afectación directa en quince departamentos, al 19 de agosto. La OPS, con datos de la UNGRD al 25 de agosto, cuenta 494 municipios en 16 departamentos.

Las vías también resultaron afectadas. Al 19 de agosto, la UNGRD contaba 449 vías afectadas, cinco aeropuertos y varios puentes. El 12 de agosto, según la IFRC, la cifra era de 125 vías dañadas y dos aeropuertos.

La energía ayuda a entender lo que pasó con los perecederos. Según el decreto, la demanda de energía no atendida en la región de Caldas, Quindío y Risaralda fue al comienzo del 82~\%, y el 13 de agosto había bajado al 5,9~\%. El aeropuerto La Nubia, en Manizales, también resultó afectado.

El decreto mide además el peso de la zona afectada en la comida del país. Según el SIPSA del DANE, los departamentos afectados aportan el 11,5~\% del abastecimiento nacional de alimentos, y el 30,8~\% en el caso de las carnes. El decreto advierte presión sobre los precios por las interrupciones del transporte. También señala que la continuidad del programa de alimentación escolar quedó sin garantía.

\subsection{Manizales y Caldas}

Las cifras de damnificados de Manizales cambiaron con la fecha y con el método de conteo. El 11 de agosto, el Presidente habló en la ciudad de unos 2.000 damnificados y 1.150 familias sin vivienda. El 13 de agosto, la Alcaldía decía haber censado a más de 2.000 personas. Los reportes de la Alcaldía y de su Unidad de Gestión del Riesgo (UGR) que recoge Asocapitales hablan de 4.000 damnificados, con cortes del 13 y el 15 de agosto. El censo que cerró el 24 de agosto registró 8.240 personas damnificadas y 3.080 predios evacuados, según el alcalde Jorge Eduardo Rojas. La misma nota informa que siete edificios con daño estructural grave serán demolidos.

Las muertes en la ciudad también tienen dos cifras. El primer balance de la Alcaldía, del 10 de agosto, contaba cuatro. Los reportes con corte al 13 y al 15 de agosto subieron la cifra a seis. Para Caldas hay dos versiones, cinco muertos según la UNGRD citada por la OPS al 25 de agosto y seis según el vicepresidente José Manuel Restrepo el 18 de agosto.

Las viviendas dañadas de Manizales aparecen con fechas que no cuadran entre los informes del equipo. Según la investigación de evidencia, Asocapitales registra 1.449 viviendas colapsadas y 1.970 averiadas con corte al 13 de agosto, y 2.707 colapsadas y 3.993 averiadas con corte al 15 de agosto. La investigación de referentes cita otro informe de Asocapitales que asigna las cifras de 2.707 y 3.993 al corte del 13 de agosto. La fecha de esas dos cifras queda \porverificar.

En Caldas, la OPS cita a la UNGRD con 27 municipios afectados, 9.511 personas afectadas, 329 heridos y tres albergues temporales con 230 personas, al 25 de agosto. Las viviendas dañadas del departamento tienen tres versiones que no coinciden. El 12 de agosto, según la IFRC, la UNGRD contaba 4.064 viviendas averiadas y 1.442 destruidas. El 18 de agosto, el vicepresidente habló de 9.400 viviendas con afectación parcial y 2.200 con afectación total. La Gobernación, con datos de solo ocho municipios, reportó el 20 de agosto 2.658 con pérdida parcial y 582 con pérdida total.

Los coliseos de la ciudad se volvieron albergue y centro de acopio. El 10 de agosto, el Coliseo Mayor y el Coliseo Menor recibieron a personas evacuadas. El albergue junto al Coliseo Menor pasó de 114 a 245 personas el 12 de agosto, y el Coliseo Mayor tenía unas 300 personas el 13 de agosto. El 24 de agosto el Coliseo Mayor seguía siendo el albergue principal. Esa cifra no encaja con la de la OPS, que el 25 de agosto contaba 230 personas en tres albergues temporales de todo Caldas. Una explicación posible, que es inferencia nuestra, es que las dos fuentes cuenten tipos distintos de albergue. El dato queda \porverificar.

La respuesta del primer día movilizó más de 324 unidades de rescate y casi 2.000 policías. El alcalde pidió no sacar vehículos a la calle para dejar paso a ambulancias y bomberos. El censo de damnificados se hizo en cuatro puntos de la ciudad, que fueron Milán, Chipre, la avenida Santander y el centro. Lo coordinaron Bomberos, Defensa Civil y Cruz Roja con la Alcaldía. Estar en ese censo era requisito para recibir el subsidio de arriendo, de 300.000 pesos mensuales por tres meses prorrogables. El comercio del centro estuvo cerrado nueve días, y la Alcaldía autorizó su reapertura gradual el 18 de agosto. Con corte al 13 de agosto, la ciudad pedía colchonetas, mantas, carpas y plantas eléctricas para el albergue.

\subsection{El corredor de los bancos de alimentos}

La ayuda alimentaria tuvo en Manizales uno de sus nodos. El 10 de agosto, la red de bancos de alimentos del país (ABACO) activó un corredor humanitario, con Armenia, Cali e Ibagué como puntos intermedios. Los bancos de Pereira y Buenaventura tuvieron que buscar alternativas por daños en sus sedes. La red anunció que daría prioridad a los albergues oficiales, los hospitales y las familias autoalbergadas.

El papel de Manizales en ese corredor tiene dos lecturas. Las dos investigaciones del equipo resumen la misma nota de Noticias RCN de forma distinta. Una dice que Manizales fue el punto principal de acopio, y la otra, que fue uno de los principales. La diferencia queda \porverificar{} abriendo la nota.

El volumen del corredor se conoce solo para toda la red. El 18 de agosto, ABACO había superado las 2.000 toneladas y focalizado a más de 9.000 familias, con 23 bancos activos. Al cierre de agosto hablaba de más de 9.470 familias en 36 municipios. Esas cifras no dicen cuánto se repartió en Manizales.

La misma red nombró sus dificultades el 18 de agosto, por boca de su gerente de Cooperación Internacional. El reto principal fue la llegada al territorio. También mencionó el transporte disponible, la articulación entre actores, el almacenamiento en los acopios y la falta de vías de acceso claras, sobre todo hacia las zonas rurales. Para entonces había tres bancos de alimentos afectados, en Buenaventura, Cartago y Pereira, y el de Pereira podría tener que demolerse.

Los perecederos quedaron en segundo plano. Según el director ejecutivo de ABACO, los cortes de energía impedían manejar alimentos que necesitan frío, y por eso la red priorizó enlatados y alimentos de larga vida. La Cruz Roja, por su parte, pidió el 14 de agosto no donar perecederos, alimentos abiertos, agua con empaque dañado ni medicamentos abiertos o por vencer. Ese día informó que había despachado más de 280 toneladas a Buenaventura, Cali, Quibdó y Pereira.

Lo que pasó en Bogotá muestra dónde se forma el cuello de botella. El centro de acopio del estadio El Campín reunió 933 toneladas en seis días. Al cierre seguía pidiendo voluntarios, montacargas y camiones para clasificar y despachar lo recibido. De esa carga, 157 toneladas le correspondían a Caldas. La investigación no registra la fecha de esa nota.

En Manizales, el Coliseo Menor funcionó como centro de acopio, coordinado por la gestora social del municipio, Juliana Londoño. El 12 de agosto recibía mercados, artículos de aseo, cobijas y colchonetas. Según ella, repartía entre más de 4.000 desplazados. Llegaron además otras ayudas, sin cifras para la ciudad. El 11 de agosto el Presidente anunció kits de aseo, alimentación, noche y cocina para Manizales. El 15 de agosto llegó una tractomula con ayuda de la Alcaldía de Bogotá. En septiembre, la prensa reportó entregas de kits de mercado de una fundación privada en Manizales, Supía y Riosucio, dentro de un total de 8.000 kits.

El abastecimiento mayorista de la ciudad bajó la semana siguiente al sismo. Según el SIPSA del DANE, la plaza Centro Galerías recibió 2.069 toneladas entre el 13 y el 19 de agosto, un 4,46~\% menos que las 2.165 de la semana anterior. Los tubérculos cayeron un 20,73~\% y las verduras un 10,02~\%. El boletín atribuye la caída nacional a problemas de recolección después del sismo, pero no vincula de forma directa el dato de Manizales con el sismo.

Las quejas documentadas hasta ahora tratan de ayudas económicas y no de comida. El 5 de octubre, la Procuraduría dijo que gobernadores y alcaldes de Risaralda, Quindío y Caldas señalaban desembolsos tardíos y una distribución desigual. El Registro Único de Damnificados cerró el 4 de septiembre y dejó por fuera a familias rurales dispersas y a familias étnicas.

\subsection{Las vías cerradas}

Dos cierres totales afectaron las salidas de la ciudad en los primeros días. El primero estaba en la vía Pereira–Chinchiná, entre los kilómetros 22 y 24. La Policía de Tránsito lo reportó el 10 de agosto, y seguía vigente el 12 de agosto a las 9:24 de la mañana y a las 7:10 de la noche. La prensa dice que el cierre afecta la conexión del sur de Caldas con Pereira. Los documentos del equipo lo cuentan como una de las salidas de Manizales. Si el tramo entre Manizales y Chinchiná estuvo cerrado queda \porverificar.

El segundo cierre estaba en la vía Manizales–Fresno, en el Alto de Letras, entre los kilómetros 36 y 40, ya en Herveo (Tolima). Bloqueaba la salida de Manizales hacia Fresno, Mariquita y Honda, y seguía vigente el 12 de agosto. Según Portafolio, la UNGRD todavía incluía los dos corredores entre los afectados en su corte del 19 de agosto.

El corredor concesionado Armenia–Pereira–Manizales tuvo derrumbes sobre la calzada en varios puntos, como Tarapacá I y II, Club Campestre–Caimo y La Paila–Río La Vieja. La Agencia Nacional de Infraestructura (ANI) lo informó el 10 de agosto. Según El Colombiano, que cita al INVÍAS, los daños típicos fueron deslizamientos traslacionales, caída de material vegetal y daños totales o parciales en las calzadas. El 25 de agosto, la OPS seguía reportando restricciones de acceso terrestre al Eje Cafetero.

Varias preguntas sobre las vías siguen sin fuente. No encontramos reportes de cierres por el sismo en la vía de La Siria ni en las salidas a Medellín y a Bogotá. La vía a Bogotá, en cambio, pasa por Letras y Fresno, que sí figuran cerradas. Tampoco encontramos la fecha de reapertura de ninguno de los dos cierres, ni reportes de embotellamientos en los accesos a la ciudad.

\subsection{Las cifras en una tabla}

La tabla~\ref{tab:cifras} reúne las cifras de esta sección con su fecha de corte, su fuente y su nivel de evidencia. Las fuentes completas, con enlace, están en el anexo.

{\footnotesize
\setlength{\tabcolsep}{4pt}
\renewcommand{\arraystretch}{1.25}
\begin{longtable}{L{0.17\linewidth} L{0.31\linewidth} L{0.13\linewidth} L{0.18\linewidth} L{0.12\linewidth}}
\caption{Cifras del sismo con su fecha de corte. Ninguna ha sido abierta todavía por una persona del equipo.}\label{tab:cifras}\\
\toprule
\textbf{Dato} & \textbf{Valor} & \textbf{Corte} & \textbf{Fuente} & \textbf{Nivel}\\
\midrule
\endfirsthead
\multicolumn{5}{l}{\small\itshape Tabla~\ref{tab:cifras}, continuación}\\
\toprule
\textbf{Dato} & \textbf{Valor} & \textbf{Corte} & \textbf{Fuente} & \textbf{Nivel}\\
\midrule
\endhead
\midrule
\multicolumn{5}{r}{\small\itshape Sigue en la página siguiente}\\
\endfoot
\bottomrule
\endlastfoot
Magnitud, hora y epicentro & 7,4 Mw; 7:34 a.\,m.; San José del Palmar (Chocó) & 10 ago 2026 & Decreto 1261, que cita al SGC & \textbf{Oficial} (alta)\\
Profundidad & 103,4 km en un considerando y unos 94 km en otro & 19 ago 2026 & Decreto 1261 & \textbf{Oficial} (media)\\
Emergencia & 30 días calendario en 15 departamentos, Caldas entre ellos & 19 ago 2026 & Decreto 1261 & \textbf{Oficial} (alta)\\
Balance nacional (reporte 36) & 314 fallecidos, 4.437 heridos y 262 desaparecidos; 262.154 damnificados en 147.464 familias & 19 ago 2026 & Decreto 1261, que cita a la UNGRD & \textbf{Oficial} (alta, preliminar)\\
Balance nacional posterior & 331 fallecidos, 4.439 heridos, 240 desaparecidos y 361.707 afectados & 25 ago 2026 & OPS, informe de situación 6, que cita a la UNGRD & \textbf{Oficial} (alta)\\
Muertos, otros cortes & 87 en capitales (10 ago); 241 (12 ago); 335 en 500 municipios (octubre) & 10 ago a oct 2026 & Asocapitales n.º 6; IFRC; Defensoría vía Minuto60 & \textbf{Prensa} (sin calificar)\\
Municipios afectados & 471 en 15 departamentos (19 ago); 494 en 16 departamentos (25 ago) & 19 y 25 ago 2026 & Decreto 1261; OPS & \textbf{Oficial} (alta)\\
Vías afectadas en el país & 449 vías, 5 aeropuertos y varios puentes (19 ago); 125 vías y 2 aeropuertos (12 ago) & 12 y 19 ago 2026 & Decreto 1261; IFRC & \textbf{Oficial} (alta); \textbf{Prensa} (sin calificar)\\
Energía en Caldas, Quindío y Risaralda & Demanda no atendida del 82~\% al comienzo y del 5,9~\% el 13 de agosto & 13 ago 2026 & Decreto 1261 & \textbf{Oficial} (alta)\\
Peso en el abastecimiento & 11,5~\% de los alimentos del país; 30,8~\% de las carnes & 19 ago 2026 & Decreto 1261, que cita al SIPSA-DANE & \textbf{Oficial} (alta)\\
Muertos en Manizales & 4 en el primer balance; 6 en los cortes del 13 y el 15 de agosto & 10, 13 y 15 ago 2026 & Noticias RCN; Asocapitales n.º 22 y n.º 27 & \textbf{Prensa} (media)\\
Muertos en Caldas & 5 (UNGRD vía OPS); 6 (vicepresidente) & 25 y 18 ago 2026 & OPS; Blu Radio & \textbf{Oficial} y \textbf{Prensa} (media)\\
Damnificados en Manizales & Unos 2.000 (11 ago); más de 2.000 censados (13 ago); 4.000 (13 y 15 ago); 8.240 personas y 3.080 predios evacuados (24 ago) & 11 a 24 ago 2026 & NTN24; Blu Radio; Asocapitales; Noticias RCN & \textbf{Prensa} (media)\\
Viviendas en Manizales & 1.449 colapsadas y 1.970 averiadas (13 ago); 2.707 y 3.993 (15 ago, otro informe las fecha el 13 ago: \porverificar) & 13 y 15 ago 2026 & Asocapitales n.º 22, n.º 24 y n.º 27 & \textbf{Prensa} (media)\\
Caldas según la UNGRD & 27 municipios; 9.511 personas afectadas; 329 heridos; 3 albergues temporales con 230 personas & 25 ago 2026 & OPS, que cita a la UNGRD & \textbf{Oficial} (alta)\\
Viviendas en Caldas & 4.064 averiadas y 1.442 destruidas (12 ago); 9.400 parciales y 2.200 totales (18 ago); 2.658 y 582 en 8 municipios (20 ago) & 12 a 20 ago 2026 & IFRC; Blu Radio (vicepresidente); Portafolio (Gobernación) & \textbf{Prensa} (media; IFRC sin calificar)\\
Albergues en Manizales & De 114 a 245 personas junto al Coliseo Menor (12 ago); unas 300 en el Coliseo Mayor (13 ago); Coliseo Mayor como albergue principal (24 ago) & 10 a 24 ago 2026 & Noticias RCN; Blu Radio & \textbf{Prensa} (media)\\
Corredor de ABACO & Más de 2.000 t y más de 9.000 familias (18 ago); más de 9.470 familias en 36 municipios (cierre de agosto) & 18 ago y 21 sep 2026 & Noticias RCN; Portafolio & \textbf{Prensa} (media)\\
Acopio de El Campín (Bogotá) & 933 t en seis días; 157 t para Caldas & Sin fecha en la investigación & Noticias RCN & \textbf{Prensa} (sin calificar)\\
Cierres de salida & Pereira–Chinchiná, km 22 a 24, y Manizales–Fresno, km 36 a 40; ambos totales y vigentes el 12 ago & 10 y 12 ago 2026 & Publimetro; Semana; Colombia.com; Minuto30 & \textbf{Prensa} (media)\\
Plaza Centro Galerías & 2.069 t entre el 13 y el 19 de agosto, 4,46~\% menos que la semana anterior & 15 a 21 ago 2026 & DANE, SIPSA, boletín semanal & \textbf{Oficial} (media)\\
\end{longtable}
}

% Sección 3
\section{Por qué en Manizales es más difícil}

\subsection{Una ciudad de ladera con pocas entradas}

Manizales está construida sobre una cuchilla, y casi todas sus vías de acceso suben o bajan por ladera. Esa descripción sale del conocimiento que el equipo tiene de su propia ciudad. No encontramos un estudio que mida cuánto depende Manizales de pocas vías de acceso, así que esa dependencia es por ahora inferencia nuestra.

Un análisis de la Universidad Javeriana de Cali sobre el sismo coincide con esa lectura. El estudio agrupó los lugares afectados en clústeres y puso a Manizales en uno que llama «urbano de ladera», con 361 puntos en el registro que usa y a 81 km del epicentro. Ese registro es abierto y no oficial.

El terreno también decide qué vehículo llega a cada sitio. Donde no entra un camión puede entrar un campero o una moto, y esa diferencia pesa en la última milla de los barrios de ladera y de las veredas. Por eso el prototipo trabaja con una flota mixta y con reglas de acceso por tipo de vehículo. Esta parte es inferencia nuestra y no viene de una medición.

\subsection{Cierres por derrumbe antes del sismo}

Los derrumbes cierran las vías de Manizales también sin terremoto. En marzo de 2017 un derrumbe cerró la vía Manizales–Fresno, que es la salida hacia Bogotá, en el sector Brisas, en el límite con Caldas. En abril de 2017, los deslizamientos en los barrios de ladera Aranjuez, Camilo Torres y Bajo Persia dejaron 17 muertos y 75 viviendas colapsadas, y obligaron a desalojar unas 400. La Cruz Roja estimó entonces que habría que reubicar unas 500 familias y revisar 11 barrios en riesgo.

Los cierres siguieron en los años siguientes. En marzo de 2021, un deslizamiento de unos 4.000~m³ por lluvias cerró por completo la doble calzada Manizales–Chinchiná en el kilómetro 11+500, y el tránsito se desvió por la vía antigua, en el sector de La Siria. En marzo de 2023, la vía hacia Mariquita se cerró de forma total e indefinida por lluvias y deslizamientos. Hubo que remover material hasta Las Letras, y el desvío fue por el Alto de La Línea. En junio de 2023, otro deslizamiento cerró la vía a Bogotá en el Alto de Letras. El INVÍAS recomendó entonces desviarse por Pereira, Armenia y La Línea, o por La Dorada y Medellín. En diciembre de 2024, un derrumbe cerca de Padua (Tolima) volvió a bloquear la vía hacia Fresno.

No todos los cierres son por derrumbe. En enero de 2018, el puente del Irra, en la vía hacia Medellín, se cerró de forma preventiva e indefinida por fallas estructurales, con desvíos por Supía, Riosucio y Anserma o por La Virginia y Pereira. Estos antecedentes son los que encontramos en prensa y en la ANI. La investigación no revisó los años 2011 y 2022, y no halló la duración de varios de estos cierres.

La ciudad tiene también capacidad propia frente al riesgo. En febrero de 2026, Manizales fue certificada como «Nodo de Resiliencia» dentro de la iniciativa MCR2030 de la UNDRR, la oficina de Naciones Unidas para la reducción del riesgo de desastres. Según Semana, la ciudad cuenta con un sistema de alerta multiamenaza y ha hecho más de 1.000 intervenciones de reducción del riesgo. La misma nota dice que no hubo muertes por deslizamientos en 2024 ni en 2025 hasta su fecha de publicación. El IDEA de la UNAL aporta además un indicador de alerta temprana, el A25, que usa la lluvia acumulada de 25 días con umbrales amarillo, naranja y rojo. Para la recuperación del territorio esta capacidad cuenta, porque un sistema de despacho puede apoyarse en ella en vez de partir de cero. Esa última idea es inferencia nuestra.

\subsection{La cadena de frío}

La cadena de frío depende de la energía. Al comienzo de la emergencia, la demanda de energía no atendida en Caldas, Quindío y Risaralda fue del 82~\%, y el 13 de agosto todavía era del 5,9~\%. El director ejecutivo de ABACO explicó que, sin energía, la red no podía manejar perecederos, y que por eso priorizó enlatados y alimentos de larga vida. La Cruz Roja pidió además no donar perecederos.

Nuestra lectura es que los perecederos quedaron casi fuera de la ayuda en los primeros días. Es inferencia nuestra, porque no encontramos ninguna cifra pública de alimentos donados que se dañaran o vencieran en Manizales. Como referencia nacional, una nota de Noticias RCN atribuye al Departamento Nacional de Planeación (DNP) una estimación de 9,76 millones de toneladas de alimentos perdidos o desperdiciados al año, el 34~\% de la oferta del país. Nadie del equipo ha abierto el documento del DNP.

\subsection{Cómo se organiza el territorio}

Manizales se divide en 12 comunas en la zona urbana y 7 corregimientos en la zona rural. Según el consolidado de juntas de acción comunal de la Alcaldía, la ciudad tiene 204 juntas, que corresponden a 131 barrios urbanos y 73 veredas. El documento no trae nombres ni coordenadas. La población proyectada para 2025 es de 472.087 habitantes, según una página que reproduce la proyección del DANE. Falta confirmar esa cifra en el propio DANE.

Las juntas de acción comunal son la forma en que barrios y veredas se organizan frente a la Alcaldía. Para DespachoX son la lista natural de quienes piden, porque la ficha del reto nombra a los comités barriales entre sus usuarios. Las juntas existen antes y después de una emergencia. Por eso un canal de pedidos montado sobre ellas podría seguir sirviendo cuando termine la fase de respuesta, que es lo que la mentoría llama recuperación del territorio. Esta idea es una hipótesis por validar con las propias juntas.

Las veredas pesan en este problema por dos razones documentadas. ABACO nombró la llegada a las zonas rurales como su mayor reto, y el Registro Único de Damnificados dejó por fuera a familias rurales dispersas. Un despacho que acepte pedidos colectivos por junta podría llegar a familias que no figuran en ese registro. Esa es una propuesta del equipo, tomada de las lecciones de otros desastres de la región.

La gente nombra los lugares a su manera. Un pedido puede decir «por la Siria» o «abajo de la Terminal», y quien despacha tiene que saber dónde queda eso. En Haití, en 2010, el conocimiento local de los nombres de lugar fue decisivo para ubicar los mensajes de ayuda. En Manizales todavía no sabemos qué comunas concentraron a los damnificados ni cuáles quedaron aisladas. Los únicos lugares que nombra la prensa son los cuatro puntos de censo de Milán, Chipre, la avenida Santander y el centro.

% Sección 4
\section{Las preguntas de la mentoría}

La mentoría dejó las cinco preguntas que siguen. Dos agentes investigan esta noche lo que falta para responderlas. Uno busca alternativas, actores y el marco oficial de recuperación. El otro reúne lo que dijeron los actores en redes, prensa e informes de gestión. Lo que depende de ellos lleva la marca \eninvestigacion.

\begin{enumerate}[leftmargin=1.6em, itemsep=0.1em]
\item ¿Por qué hacemos esto y para qué sirve?
\item ¿A quién le sirve y a quién impacta?
\item ¿Hay algo parecido? ¿Por qué no se usa en Manizales o en la región?
\item ¿Qué prototipo va a medir el efecto en la región?
\item ¿Cómo es el piloto de 30 a 90 días para la recuperación del territorio?
\end{enumerate}

\subsection{Por qué hacemos esto y para qué sirve}

\paragraph{Lo que sabemos} La comida llegó a Manizales, y la dificultad estuvo en repartirla. ABACO movió más de 2.000 toneladas en la primera semana y nombró la llegada al territorio como su reto principal. En Bogotá, el acopio de El Campín reunió carga más rápido de lo que podía despachar. Los cortes de energía dejaron por fuera a los perecederos, y dos vías de salida seguían cerradas dos días después del sismo.

DespachoX sirve para responder cada mañana la pregunta de quien despacha. Con lo que hay en bodega y los vehículos disponibles, esa persona tiene que decidir qué le manda a quién, en qué orden y por dónde. El resultado que buscamos es que, con la misma comida y los mismos vehículos, más pedidos urgentes lleguen a tiempo y menos alimentos se dañen. Esos son los dos indicadores del reto.

La trazabilidad responde a otra necesidad, la confianza. En Colombia hay antecedentes de ayudas que se perdieron o se desviaron, como los 80 carrotanques que la UNGRD compró para La Guajira y que terminaron deteriorados sin uso. En 2026 también hubo preguntas públicas sobre el destino de las donaciones. Un código por paquete, confirmado al entregarse, permite mostrar a quien donó que su ayuda llegó.

\paragraph{Lo que sigue abierto} La ficha del reto dice que el desabastecimiento llega «en cuestión de 48 horas», y no encontramos fuente que lo respalde para este sismo. Por eso no lo usamos como hecho. Tampoco sabemos si en Manizales se vencieron alimentos donados, ni cuántas familias recibieron ayuda alimentaria en la ciudad. La pregunta de fondo, cómo despacha hoy el Banco de Alimentos de Manizales, está \pendiente.

\subsection{A quién le sirve y a quién impacta}

\paragraph{Lo que sabemos} La ficha del reto nombra como usuarios a la Cruz Roja, la Defensa Civil, el Banco de Alimentos, los transportadores y los comités barriales. El equipo pone en el centro a la persona que despacha. Puede ser una funcionaria del Banco de Alimentos con años de experiencia o un voluntario que llegó ayer, y en ambos casos decide con información incompleta. Le siguen el líder de barrio, que pide en nombre de sus vecinos, y el conductor, que conoce el estado real de las vías mejor que cualquier mapa. Las familias no van a usar la herramienta. Lo que notan es si la comida llegó a tiempo y en buen estado.

La evidencia señala a tres grupos de afectados en particular. Están las personas en albergues como el Coliseo Mayor y las familias autoalbergadas, que ABACO puso entre sus prioridades. Están también las familias rurales dispersas que quedaron fuera del Registro Único de Damnificados.

El sistema tiene un costo que ya medimos en simulación. Al atender primero lo crítico, el punto peor atendido espera más que con el despacho por orden de llegada. Por eso el sistema propone y una persona aprueba.

\paragraph{Lo que sigue abierto} Qué actores operaron de verdad en Manizales y qué dijeron está \eninvestigacion. Qué comunas concentraron a los damnificados sigue \pendiente. Si el Banco de Alimentos de Manizales y la Cruz Roja de Caldas se reconocen como usuarios también está \pendiente.

\subsection{Si hay algo parecido y por qué no se usa aquí}

\paragraph{Lo que sabemos} Existen piezas sueltas. Hay herramientas abiertas que calculan rutas con vías bloqueadas, como el openrouteservice para desastres de HeiGIT, VROOM y Valhalla. Otras registran entregas, como LMMS de World Vision, SCOPE del Programa Mundial de Alimentos (PMA) y KoBo, y otras recogen reportes ciudadanos, como Ushahidi. La OPS creó en los años noventa un sistema regional de manejo de suministros humanitarios, SUMA/LSS, que Colombia incorporó en su normativa.

Después del sismo aparecieron plataformas para informar y para recaudar. Cuidar a Colombia, de la Universidad EAFIT, publicó un directorio verificado con 58 puntos de acopio en 32 ciudades. Un bot de WhatsApp de JumpCube con AWS reunió información oficial para los damnificados. ABACO recibió donaciones en línea y Rappi habilitó un botón para donar.

La región tiene además antecedentes útiles. En Haití, en 2010, Mission 4636 convertía mensajes de texto en reportes estructurados y ubicados en el mapa en unos cinco minutos. En Brasil, en 2024, la web Salva RS! se construyó en unas 48 horas para ordenar los pedidos de rescate que circulaban por WhatsApp, y la Defensa Civil la usó desde el tercer día. En México, en 2017, Verificado19S comprobaba cada reporte y publicaba necesidades concretas por zona.

Lo que no encontramos es una herramienta abierta que una todo el recorrido. Ese recorrido va del mensaje de un vecino a la ruta de un vehículo, con el cierre que reportó un conductor y el vencimiento de lo que va en la carga. Las plataformas de 2026 que encontramos se dedican a informar o a recaudar, y ninguna de ellas despacha pedidos.

Nuestra explicación de por qué algo así no se usa en Manizales es que no existe unido. Cada pieza resuelve una parte, y juntarlas bajo presión supera lo que un equipo de emergencia puede montar en pocos días. Esta explicación es inferencia nuestra. De SUMA/LSS no sabemos si alguien lo usa en Caldas, porque la investigación no verificó su uso por parte de la comunidad.

\paragraph{Lo que sigue abierto} Qué alternativas usan hoy la Alcaldía, la UGR y los organismos de socorro está \eninvestigacion. Si el Banco de Alimentos de Manizales usa algún sistema propio está \pendiente.

\subsection{Qué prototipo va a medir el efecto en la región}

\paragraph{Lo que sabemos} El prototipo es una mesa de despacho. Los comités barriales piden por chat, con texto o con voz, y los conductores reportan las vías bloqueadas por el mismo canal. Un modelo de lenguaje lee y estructura cada mensaje. Luego un optimizador asigna los pedidos a una flota mixta y calcula las rutas con lo que hay en bodega y lo que vence primero. El despachador aprueba cada plan, y cada paquete lleva un código QR hasta que alguien confirma la entrega. La regla del equipo es que el modelo lee, el optimizador calcula y la persona decide.

El 20 de noviembre el sistema corrió por primera vez de punta a punta en la nube. Un mensaje recibido por el bot se convirtió en un pedido confirmado y en un plan calculado. En esa prueba nadie aprobó el plan ni probó un cierre.

Hoy medimos los dos indicadores del reto solo en simulación, y la sección~\ref{sec:prototipo} da las cifras. En un piloto, la confirmación de cada entrega con el código QR daría los datos para medirlos en la operación real. Esto último es inferencia nuestra.

\paragraph{Lo que sigue abierto} No hay todavía ninguna medición con datos reales de Manizales. La simulación usa supuestos nuestros sobre demanda, flota, inventario y cierres, y no mide el valor de recibir los pedidos por chat ni el de la trazabilidad.

\subsection{Cómo es el piloto de 30 a 90 días}

\paragraph{Lo que sabemos} Tenemos una lista de aliados posibles y lo que se le pediría a cada uno. ABACO y el Banco de Alimentos de Manizales serían los usuarios del piloto y podrían aportar datos de los despachos de agosto. A la UGR de Manizales se le pediría la lista de albergues y de vías cerradas, y el aval para operar. La Cruz Roja de Caldas podría probar el tablero en el reparto. El equipo humanitario de OpenStreetMap (HOT) reporta una ortofoto de dron de Manizales del 19 de agosto, y el IDEA de la UNAL tiene datos de lluvia. Exportar la trazabilidad en un formato compatible con SUMA/LSS sería un buen hito para el día 60.

Del marco oficial conocemos una pieza. El Fondo Milagro, creado por el Decreto Legislativo 1379 del 9 de septiembre, financia la reconstrucción. Lo administra Fiducoldex y recibió su primer billón de pesos el 7 de octubre. Su objeto son la demolición, la vivienda y los colegios, y la ayuda alimentaria queda por fuera.

\paragraph{Lo que sigue abierto} El diseño del piloto se escribe esta noche (sección~\ref{sec:piloto}). El marco oficial de recuperación, más allá del Fondo Milagro, está \eninvestigacion. Ningún aliado ha confirmado que nos recibiría, así que ese punto está \pendiente.

% Sección 5
\section{Lo que dicen los actores}

Las fuentes públicas nombran a pocos actores con nombre propio. En la Alcaldía de Manizales hablaron el alcalde, Jorge Eduardo Rojas, y la gestora social, Juliana Londoño, que coordinó el acopio del Coliseo Menor. En la Gobernación de Caldas aparece el gobernador, Henry Gutiérrez. En ABACO hablaron el director ejecutivo, Juan Carlos Buitrago, y la gerente de Cooperación Internacional, Norma Fernanda Alonso. El censo de damnificados lo coordinaron Bomberos, Defensa Civil y Cruz Roja con la Alcaldía, y lo acompañaron la Personería, la Defensoría, el ICBF y la fuerza pública.

Otros actores clave no tienen declaraciones públicas que hayamos encontrado. Es el caso del Banco de Alimentos de Manizales, que la investigación asocia con la Arquidiócesis sin haberlo confirmado. Tampoco encontramos declaraciones de la Cruz Roja de Caldas, la Defensa Civil de Caldas, la UDEGER ni la UGR de Manizales. La búsqueda de lo que dijeron los actores en redes, prensa e informes de gestión está \eninvestigacion.

La tabla~\ref{tab:entrevista} recoge las ocho preguntas que el equipo quiere hacer en entrevistas a quienes despacharon, pidieron o transportaron comida después del sismo. Ninguna respuesta viene todavía de una entrevista, porque el equipo no ha hablado con el Banco de Alimentos de Manizales ni con los organismos de socorro de Caldas. Lo que sabemos sale de prensa e informes públicos.

{\footnotesize
\setlength{\tabcolsep}{4pt}
\renewcommand{\arraystretch}{1.3}
\begin{longtable}{L{0.025\linewidth} L{0.17\linewidth} L{0.50\linewidth} L{0.22\linewidth}}
\caption{Las ocho preguntas de entrevista, con lo que ya sabemos y su fuente.}\label{tab:entrevista}\\
\toprule
\textbf{\#} & \textbf{Pregunta} & \textbf{Lo que ya sabemos} & \textbf{Fuente}\\
\midrule
\endfirsthead
\multicolumn{4}{l}{\small\itshape Tabla~\ref{tab:entrevista}, continuación}\\
\toprule
\textbf{\#} & \textbf{Pregunta} & \textbf{Lo que ya sabemos} & \textbf{Fuente}\\
\midrule
\endhead
\midrule
\multicolumn{4}{r}{\small\itshape Sigue en la página siguiente}\\
\endfoot
\bottomrule
\endlastfoot
1 & ¿Cómo decidían el despacho? & ABACO anunció que daría prioridad a los albergues oficiales, los hospitales y las familias autoalbergadas. Según la investigación de referentes, ABACO registra fechas de vencimiento; la fuente directa de ese dato está \porverificar. Cómo decide paso a paso el Banco de Alimentos de Manizales está \pendiente. & Noticias RCN, 10 y 11 ago 2026 (director ejecutivo de ABACO)\\
2 & ¿Cómo llegaban los pedidos? & \pendiente. No hay información pública sobre cómo pedían los comités barriales en Manizales. Que los pedidos lleguen desordenados por WhatsApp y por teléfono es una hipótesis del equipo. Como referencia, en Rio Grande do Sul (Brasil, 2024) los pedidos de rescate circularon por grupos de WhatsApp y se repetían. & Fast Company Brasil (referente regional)\\
3 & ¿Cómo se enteraban de los cierres? & La prensa citó reportes de la Policía de Tránsito, y la ANI informó el estado de sus corredores. No encontramos comunicados directos del INVÍAS, la Gobernación ni la Alcaldía. Que el reporte del conductor viaje por teléfono, tarde y sin registro es una hipótesis del equipo. Cómo se enteraban quienes despachaban está \pendiente. & Publimetro, Semana y Colombia.com (10 y 12 ago 2026); ANI (10 ago 2026)\\
4 & ¿Qué comida se dañó? & ABACO explicó que sin energía no podía manejar perecederos y que priorizó enlatados y alimentos de larga vida. La Cruz Roja pidió no donar perecederos. La cifra de alimentos donados que se dañaron o vencieron en Manizales está \pendiente. & La República (entrevista, sin fecha en la investigación); Infobae, 14 ago 2026\\
5 & ¿Qué plataformas usaron? & En el país circularon Cuidar a Colombia (EAFIT), un bot de WhatsApp de JumpCube con AWS y los botones de donación de ABACO y Rappi. El Estado usó el Registro Único de Damnificados, que cerró el 4 de septiembre. Cuáles usaron los actores de Manizales está \pendiente. & El Colombiano; Noticias RCN; Infobae, 11 ago 2026; Vanguardia, 5 oct 2026\\
6 & ¿Qué les haría falta? & ABACO nombró la llegada al territorio, el transporte disponible, la articulación entre actores, el almacenamiento en los acopios y vías de acceso claras hacia las zonas rurales. El acopio de El Campín pidió camiones, montacargas y voluntarios. Con corte al 13 de agosto, Manizales pedía colchonetas, mantas, carpas y plantas eléctricas para el albergue. Lo que diría el Banco de Alimentos de Manizales está \pendiente. & Noticias RCN, 18 ago 2026; Noticias RCN (El Campín, sin fecha); Asocapitales n.º 24\\
7 & ¿Quién usaría la herramienta? & La ficha del reto nombra a la Cruz Roja, la Defensa Civil, el Banco de Alimentos, los transportadores y los comités barriales. El equipo propone a ABACO y al Banco de Alimentos como usuario principal. La confirmación de los propios actores está \pendiente. & Ficha de la vertical 04, citada en \ruta{docs/PROBLEMA.md}; \ruta{docs/ESTADO_DEL_ARTE.md}\\
8 & ¿Quién nos recibiría para un piloto? & Los candidatos son ABACO, el Banco de Alimentos de Manizales, la UGR de Manizales, la Cruz Roja de Caldas, el PMA, el equipo humanitario de OpenStreetMap (HOT), el IDEA de la UNAL y Asocapitales. El nombre y el canal oficial del Banco de Alimentos de Manizales no están confirmados. Ninguna respuesta todavía, así que está \pendiente. & \ruta{docs/investigacion/05_datos_abiertos_y_aliados.md}, parte C\\
\end{longtable}
}

% Sección 6
\section{Qué mide hoy el prototipo}\label{sec:prototipo}

\subsection{Los indicadores del reto}

La ficha de la vertical 04 mide el resultado con dos indicadores. El primero es el porcentaje de entregas a tiempo en puntos críticos, como un albergue o un comedor comunitario. El segundo es el desperdicio de alimentos perecederos, que debe quedar por debajo del 5~\%. Hoy los medimos solo en simulación, porque el prototipo no ha operado con datos reales de Manizales.

\subsection{Cómo se simuló}

La simulación compara cuatro formas de despachar la misma jornada, con los mismos pedidos, la misma flota y los mismos cierres. Usa el grafo de conducción de Manizales y Villamaría de OpenStreetMap. Cada escenario tiene entre 30 y 80 puntos de entrega, de los cuales el 20~\% son críticos. La flota mezcla camiones, camperos y motos, y en cada jornada se cierran entre 3 y 8 tramos de vía. Nadie conoce un cierre hasta que un vehículo intenta pasar por él. La jornada dura 14 horas, y lo que no se entregó para entonces cuenta como no entregado.

Las cuatro políticas se distinguen por cómo deciden. La línea base, llamada ingenua, atiende por orden de llegada y saca los lotes de bodega por orden de ingreso. La política experta imita a un despachador con oficio, que atiende primero lo urgente y saca primero lo que vence. Las otras dos usan el optimizador de DespachoX. Una lo llama una sola vez al inicio, y la otra, que es DespachoX completo, recalcula con cada cierre, cada tanda de pedidos y cada regreso de un vehículo al acopio.

Cada nivel de carga se corrió con 30 semillas, es decir, 30 jornadas sintéticas, las mismas para todas las políticas. Los resultados se comparan por pares, semilla a semilla, con un intervalo de confianza del 95~\%. Todos los datos de demanda, flota, inventario y cierres son sintéticos. Las velocidades, los tiempos y los plazos son supuestos del equipo, y ninguno viene de una operación real del Banco de Alimentos ni de los organismos de socorro. Nadie del equipo ha revisado todavía estos resultados.

\subsection{Tres cifras con su línea base}

La tabla~\ref{tab:simulacion} trae las tres cifras que el agente de simulación propuso para el pitch, con su línea base y su número de escenarios. Todas son simulación con supuestos nuestros, y su nivel de evidencia es \textbf{Inferencia}. Las fuentes son las entradas de las 4:57~p.\,m. del canal del equipo (\ruta{torre/CANAL.md}) y la sección 8 de \ruta{docs/SIMULACION.md}, en la rama principal del repositorio. La corrida con flota corta todavía no está descrita en \ruta{docs/SIMULACION.md}, así que su número de semillas sale de su archivo de resultados.

{\footnotesize
\setlength{\tabcolsep}{4pt}
\renewcommand{\arraystretch}{1.3}
\begin{longtable}{L{0.20\linewidth} L{0.15\linewidth} L{0.24\linewidth} L{0.32\linewidth}}
\caption{Cifras de simulación con su línea base. Ninguna es una medición en la operación real.}\label{tab:simulacion}\\
\toprule
\textbf{Indicador} & \textbf{DespachoX} & \textbf{Línea base} & \textbf{Escenarios y condiciones}\\
\midrule
\endfirsthead
\multicolumn{4}{l}{\small\itshape Tabla~\ref{tab:simulacion}, continuación}\\
\toprule
\textbf{Indicador} & \textbf{DespachoX} & \textbf{Línea base} & \textbf{Escenarios y condiciones}\\
\midrule
\endhead
\midrule
\multicolumn{4}{r}{\small\itshape Sigue en la página siguiente}\\
\endfoot
\bottomrule
\endlastfoot
(1) Entregas a tiempo en puntos críticos, carga media & 99,5~\% & 70,9~\% con despacho por orden de llegada. Diferencia de +28,6 puntos, con intervalo [22,4; 34,7] & Simulación con 30 semillas nuevas que no intervinieron en ninguna decisión. Usa un peso de 3 para los puntos críticos (\texttt{FACTOR\_CRITICO}), propuesto a las 4:57~p.\,m. y que a esa hora no era todavía el valor por defecto del código.\\
(2) Perecederos vencidos, carga media & 0,0~\% & 1,8~\% al sacar de bodega por orden de ingreso & Simulación con 30 semillas. Con dos camperos menos, también en 30 semillas, el orden de ingreso llega al 5,7~\% y pasa la meta del 5~\%, mientras DespachoX sigue en 0,0~\%.\\
(3) Entregas a tiempo en todos los puntos, con un 10~\% de pedidos mal ubicados & +6,0 puntos, con intervalo [4,5; 7,5] & El mismo sistema sin el paso de confirmar la ubicación con quien pidió & Simulación con 30 semillas, carga media. El 10~\% es un supuesto sin fuente, porque no sabemos cuántos pedidos llegan mal ubicados.\\
\end{longtable}
}

La primera cifra cambia poco con ese peso. Con el peso de 1 que usa hoy el código, la primera ronda dio una diferencia parecida en carga media, de 28,9 puntos con intervalo [23,2; 34,7]. La ventaja frente al despacho por orden de llegada crece con la carga. En la primera ronda fue de 12, 29 y 46 puntos en carga holgada, media y alta, y no sabemos cuál de esas cargas se parece a la real.

\subsection{Lo que hay que saber decir frente a un despachador experto}

Un despachador con experiencia consigue casi lo mismo que DespachoX en los dos indicadores del reto, y conviene decirlo antes de que lo diga el jurado. En entregas a tiempo en puntos críticos, DespachoX empata con la política experta. En las 30 semillas nuevas, con el peso de 3, la diferencia fue de 0,3, 1,7 y 0,8 puntos según la carga, y todos los intervalos incluyen el cero. En entregas a tiempo de todos los puntos, DespachoX queda unos 2 puntos por debajo en carga media y alta. En perecederos, los dos quedan en 0,0~\%, porque el desperdicio lo resuelve una regla de bodega y no el ruteo.

Hay dos situaciones simuladas en que DespachoX supera a la experta. Con la flota corta, de dos camperos menos, en carga media y con 30 semillas, DespachoX entrega a tiempo 3,4 puntos más en todos los puntos, con intervalo [1,3; 5,5]. Con un 10~\% de pedidos mal ubicados, también en carga media y con 30 semillas, confirmar la ubicación le da 3,0 puntos más en puntos críticos [0,9; 5,2] y 3,4 más en todos los puntos [2,1; 4,8]. Esta segunda ventaja depende de un supuesto, que el despachador experto no llama a confirmar cada pedido. Si lo hiciera, tendría el mismo beneficio a costa de su tiempo, que la simulación no modela.

DespachoX también recorre menos kilómetros. En las semillas nuevas, su flota hizo entre 47 y 60 km menos por jornada que la política experta. En la primera ronda la diferencia equivalía a entre el 5~\% y el 11~\% del recorrido.

La prioridad tiene un costo en equidad. Atender primero lo crítico deja para el final a los puntos de baja urgencia. En la primera ronda, con DespachoX el punto peor atendido esperó 2,3 horas más en carga media y 4,3 horas más en carga alta que con el despacho por orden de llegada.

Los cierres viales pesan poco en estos resultados. Con 0, 5 o 10 cierres por jornada, la ventaja sobre la línea base casi no cambia, porque viene de priorizar los pedidos críticos. La red urbana simulada tiene caminos alternos, y el aviso por radio desvía a todos los conductores por igual. Esto puede cambiar en la red de la región, donde hay menos alternativas. Esa última idea es una hipótesis que la simulación todavía no prueba.

Frente a un experto, entonces, el argumento de DespachoX es la disponibilidad de su criterio. El sistema pone las reglas de un despachador con oficio al alcance de cualquier voluntario, a cualquier hora. Además las combina con información que hoy llega tarde o no llega, como los pedidos del chat, los cierres que reportan los conductores y la confirmación de cada entrega. La simulación todavía no mide el valor de esa información, salvo en el experimento de la ubicación.

\subsection{Lo que la simulación no prueba}

\begin{itemize}
\item No prueba que DespachoX supere a un buen despachador humano en los indicadores del reto.
\item No mide el valor de recibir los pedidos por chat ni el de la trazabilidad, salvo el paso de confirmar la ubicación.
\item No valida sus supuestos. Las velocidades, los tiempos de carga y de entrega, los plazos, la flota, la demanda y los cierres son nuestros.
\item No representa bien los perecederos. Hay un solo acopio, tanto inventario perecedero como demanda, y lotes que vencen entre 6 y 48 horas.
\item No incluye errores humanos ni del modelo de lenguaje, y cubre una sola jornada de 14 horas.
\end{itemize}

Según el canal, a las 4:57~p.\,m. seguían en curso dos corridas. Una usa 3 segundos de búsqueda por llamada al optimizador, y la otra corre la carga media sobre la red ampliada de la región. El agente de simulación esperaba tenerlas hacia las 5:45~p.\,m.

% Sección 7
\section{El piloto de 30 a 90 días}\label{sec:piloto}

En discusión entre el capitán y el coordinador; se escribe esta noche.

% Sección 8
\section{Lo que no sabemos todavía y a quién preguntarlo}

Las preguntas abiertas se responden mejor con una persona que estuvo en la operación que con más prensa. Una respuesta de alguien de gestión del riesgo o de un organismo de socorro, citada con nombre y cargo, pesa más que cualquier cifra de este informe. La tabla~\ref{tab:vacios} pone primero los vacíos que más afectan al despacho. La columna de la derecha propone a quién preguntar; cuando la investigación no lo dice, la propuesta es nuestra.

{\footnotesize
\setlength{\tabcolsep}{4pt}
\renewcommand{\arraystretch}{1.3}
\begin{longtable}{L{0.025\linewidth} L{0.35\linewidth} L{0.28\linewidth} L{0.26\linewidth}}
\caption{Vacíos de información y a quién preguntar.}\label{tab:vacios}\\
\toprule
\textbf{\#} & \textbf{Lo que no sabemos} & \textbf{Por qué importa} & \textbf{A quién preguntarlo}\\
\midrule
\endfirsthead
\multicolumn{4}{l}{\small\itshape Tabla~\ref{tab:vacios}, continuación}\\
\toprule
\textbf{\#} & \textbf{Lo que no sabemos} & \textbf{Por qué importa} & \textbf{A quién preguntarlo}\\
\midrule
\endhead
\midrule
\multicolumn{4}{r}{\small\itshape Sigue en la página siguiente}\\
\endfoot
\bottomrule
\endlastfoot
1 & Cómo despacha hoy, paso a paso, el Banco de Alimentos de Manizales, quién lo dirige y qué inventario y comités atiende & Es el usuario principal que propone el equipo, y nadie ha hablado con él & Banco de Alimentos de Manizales; ABACO\\
2 & Qué vías de acceso se cerraron por el sismo, desde cuándo y hasta cuándo, incluidas La Siria y las salidas a Medellín y a Bogotá & Solo hay prensa de los dos primeros días, y sin la duración no sabemos cuánto pesaron los cierres & UGR de Manizales; Policía de Tránsito; INVÍAS; Gobernación de Caldas\\
3 & Si hubo desabastecimiento de alimentos en las primeras 48 horas, y en qué barrios & La ficha del reto lo afirma y no encontramos fuente & Alguien de gestión del riesgo presente en la hackathon; UGR de Manizales\\
4 & Si se dañaron o vencieron alimentos donados en Manizales, y cuántos & Es el segundo indicador del reto & ABACO; Banco de Alimentos de Manizales; Cruz Roja de Caldas\\
5 & Cuántas raciones, kits y familias se atendieron en Manizales, y desde qué centros de acopio & Solo hay cifras nacionales de ABACO y de donantes privados & Gestora social de Manizales (acopio del Coliseo Menor); ABACO\\
6 & Qué comunas y barrios concentraron a los damnificados y cuáles quedaron aislados & Define dónde están los puntos críticos & UGR de Manizales\\
7 & Cuántas familias hubo en albergues y cómo cambió esa cifra semana a semana & Hay cifras de personas y no de familias & UGR de Manizales; Alcaldía\\
8 & Cómo operan los comités barriales y los transportadores, y qué datos tienen & Son quienes piden y quienes llevan, y no hay información pública & Juntas de acción comunal; transportadores; Alcaldía (consolidado de juntas)\\
9 & Si hubo embotellamientos críticos en los accesos a la ciudad & La ficha los menciona y solo hay antecedentes de 2021 a 2024 & Policía de Tránsito; UGR de Manizales\\
10 & Qué dicen la Cruz Roja de Caldas, la Defensa Civil de Caldas, la UDEGER y la UGR de Manizales, y quiénes las dirigen & No encontramos declaraciones propias & Cada entidad, de forma directa\\
11 & Cuántos pedidos llegan mal ubicados & La simulación supone un 10~\% sin fuente & Comités barriales; Banco de Alimentos de Manizales\\
12 & Cuántos pedidos tiene una jornada real & La ventaja simulada va de 12 a 46 puntos según la carga & Banco de Alimentos de Manizales (datos de despachos de agosto)\\
13 & Si existen ERIS-UNAL, el software de ruteo adaptativo y los datos de la ficha 04 & El banco de capacidades del sitio no tiene rastro público de ellos & Mesa de la hackathon\\
14 & Si la formulación del ruteo y el papel del despachador son razonables & Da respaldo académico local & William Sarache, Diana Cárdenas o Juan Camilo López Vargas, investigadores de la Sede Manizales en logística humanitaria\\
15 & Cómo cambiaron los precios de los alimentos en Manizales & Solo hay datos nacionales y el boletín de Centro Galerías & DANE (SIPSA); Fenalco Caldas\\
16 & La población oficial de Manizales & La cifra viene de una página que reproduce al DANE & DANE\\
17 & Qué alternativas usan hoy las entidades y qué marco oficial rige la recuperación & Responde a dos preguntas de la mentoría & \eninvestigacion\\
18 & Qué dijeron los actores en redes, prensa e informes de gestión & Sustituye en parte a las entrevistas que faltan & \eninvestigacion\\
\end{longtable}
}

% Anexo: tabla de fuentes
\clearpage
\section*{Anexo. Tabla de fuentes}
\addcontentsline{toc}{section}{Anexo. Tabla de fuentes}

La tabla lista las fuentes externas que respaldan los datos de este informe y los documentos internos de donde salen. Las fechas de publicación o de corte son las que registran los informes de investigación del equipo. Cuando esos informes no registran la fecha, la tabla lo dice. La fecha de consulta es la de esos informes, que prepararon agentes de investigación el 20 de noviembre de 2026. Ninguna fuente ha sido abierta todavía por una persona del equipo. Quien abra una debe cambiar el «no» por su nombre y la hora, y anotarla en \ruta{docs/FUENTES.md}.

{\footnotesize
\setlength{\tabcolsep}{4pt}
\renewcommand{\arraystretch}{1.25}
\begin{longtable}{L{0.35\linewidth} L{0.27\linewidth} L{0.12\linewidth} L{0.085\linewidth} L{0.085\linewidth}}
\caption{Fuentes del informe. Ninguna ha sido abierta todavía por una persona del equipo.}\label{tab:fuentes}\\
\toprule
\textbf{Fuente} & \textbf{Qué respalda} & \textbf{Publicación o corte} & \textbf{Consulta} & \textbf{Abierta por una persona}\\
\midrule
\endfirsthead
\multicolumn{5}{l}{\small\itshape Tabla~\ref{tab:fuentes}, continuación}\\
\toprule
\textbf{Fuente} & \textbf{Qué respalda} & \textbf{Publicación o corte} & \textbf{Consulta} & \textbf{Abierta por una persona}\\
\midrule
\endhead
\midrule
\multicolumn{5}{r}{\small\itshape Sigue en la página siguiente}\\
\endfoot
\bottomrule
\endlastfoot
\multicolumn{5}{l}{\textbf{El sismo y la emergencia}}\\
Decreto 1261 de 2026 (Presidencia), texto en PDF de Camacol.\newline \url{https://camacol.co/sites/default/files/descargables/DECRETO%201261de%202026_EMERGENCIA_compressed.pdf} & Magnitud, hora, epicentro y profundidad; emergencia; balance de la UNGRD al 19 de agosto; vías afectadas; energía; peso en el abastecimiento; alimentación escolar & 19 ago 2026 & 20 nov 2026 & no\\
OPS/PAHO, informe de situación 6.\newline \url{https://www.paho.org/sites/default/files/2026/08/sitrep6-colombia-sismo-26082026-es.pdf} & Balance de la UNGRD al 25 de agosto; municipios afectados; cifras de Caldas; acceso terrestre al Eje Cafetero & 25 ago 2026 & 20 nov 2026 & no\\
Asocapitales, informe n.º 6.\newline \url{https://www.asocapitales.co/sites/default/files/2026-08/informe-ejecutivo-sismo-nacional-10-08-2026.pdf} & 87 muertos en las capitales & 10 ago 2026 & 20 nov 2026 & no\\
IFRC, solicitud DREF MDRCO036.\newline \url{https://go-api.ifrc.org/api/DownloadFile/96426/MDRCO036do} & Balance de la UNGRD al 12 de agosto; viviendas de Caldas; puesto de mando unificado; Cruz Roja & 12 ago 2026 & 20 nov 2026 & no\\
Minuto60, «dos meses después».\newline \url{https://minuto60.com/colombia/terremoto-colombia-dos-meses-despues-como-va-reconstruccion/16930} & Cifra de la Defensoría (335 muertos); Fondo Milagro & Octubre de 2026, sin día en la investigación & 20 nov 2026 & no\\
Portafolio, servicios tras el terremoto.\newline \url{https://www.portafolio.co/economia/regiones/servicios-tras-el-terremoto-energia-llega-al-99-pero-siguen-restricciones-en-vias-y-aeropuertos-500954} & Corredores afectados según la UNGRD & Agosto de 2026, corte del 19 ago & 20 nov 2026 & no\\
\midrule
\multicolumn{5}{l}{\textbf{Manizales y Caldas}}\\
Asocapitales, informe n.º 22.\newline \url{https://www.asocapitales.co/sites/default/files/2026-08/informe-terremoto-capitales-no-22.pdf} & Muertos, viviendas y damnificados de Manizales (reporte de la Alcaldía y la UGR) & 13 ago 2026 & 20 nov 2026 & no\\
Asocapitales, informe n.º 24.\newline \url{https://www.asocapitales.co/sites/default/files/2026-08/informe-terremoto-capitales-no-24.pdf} & Heridos, viviendas y necesidades del albergue en Manizales & 13 ago 2026, según la investigación de referentes (\porverificar) & 20 nov 2026 & no\\
Asocapitales, informe n.º 27.\newline \url{https://www.asocapitales.co/sites/default/files/2026-08/informe-terremoto-capitales-no-27.pdf} & Muertos, 4.000 damnificados, viviendas e instituciones educativas de Manizales & 15 ago 2026 & 20 nov 2026 & no\\
Noticias RCN, balance de daños en Manizales.\newline \url{https://www.noticiasrcn.com/colombia/balance-de-danos-y-afectaciones-en-manizales-tras-fuerte-temblor-1048889} & Primer balance (4 muertos); respuesta del primer día & 10 ago 2026 & 20 nov 2026 & no\\
Noticias RCN, demolición de siete edificios.\newline \url{https://www.noticiasrcn.com/colombia/manizales-prepara-la-demolicion-de-7-edificios-que-presentan-danos-estructurales-1053593} & Censo de 8.240 damnificados y 3.080 predios evacuados & 24 ago 2026 & 20 nov 2026 & no\\
Noticias RCN, el Coliseo Menor como centro de acopio.\newline \url{https://www.noticiasrcn.com/colombia/manizales-convierte-al-coliseo-menor-en-centro-de-acopio-tras-el-terremoto-1049719} & Albergues en los coliseos; acopio del Coliseo Menor & 12 ago 2026 & 20 nov 2026 & no\\
Blu Radio, censo de damnificados en Manizales.\newline \url{https://www.bluradio.com/nacion/manizales-censa-a-mas-de-2-000-damnificados-y-anuncia-subsidios-de-arriendo-desde-la-proxima-semana-rg10} & Más de 2.000 censados; puntos de censo; subsidio de arriendo; coordinación del censo; Coliseo Mayor & 13 ago 2026 & 20 nov 2026 & no\\
NTN24, el Presidente en Manizales.\newline \url{https://www.ntn24.com/noticias-actualidad/desde-manizales-abelardo-de-la-espriella-anuncio-medidas-para-los-damnificados-de-caldas-tras-el-terremoto-la-necesidad-mas-urgente-en-este-departamento-es-el-tema-de-vivienda-641072} & Unos 2.000 damnificados y 1.150 familias sin vivienda; kits del Gobierno & 11 ago 2026 & 20 nov 2026 & no\\
Blu Radio, el vicepresidente sobre Caldas.\newline \url{https://www.bluradio.com/nacion/reconstruccion-de-caldas-luego-de-terremoto-cuesta-1-4-billones-jose-manuel-restrepo-so35} & Muertos y viviendas de Caldas según el vicepresidente & 18 ago 2026 & 20 nov 2026 & no\\
Portafolio, la Gobernación de Caldas.\newline \url{https://www.portafolio.co/economia/regiones/gobernacion-de-caldas-pedira-credito-por-80-000-millones-para-reconstruir-viviendas-tras-terremoto-500843} & Viviendas en ocho municipios; nombre del gobernador & 20 ago 2026 & 20 nov 2026 & no\\
Noticias RCN, reapertura del comercio.\newline \url{https://www.noticiasrcn.com/colombia/manizales-reactiva-su-comercio-tras-nueve-dias-cerrado-por-terremoto-1051241} & Nueve días de cierre del comercio del centro & 18 ago 2026 & 20 nov 2026 & no\\
Vanguardia, la Procuraduría pide agilizar ayudas.\newline \url{https://www.vanguardia.com/colombia/2026/10/05/procuraduria-pide-a-la-ungrd-agilizar-ayudas-para-damnificados-por-el-terremoto/} & Quejas por ayudas económicas; cierre del Registro Único de Damnificados & 5 oct 2026 & 20 nov 2026 & no\\
\midrule
\multicolumn{5}{l}{\textbf{Ayuda alimentaria y abastecimiento}}\\
Noticias RCN, corredor humanitario de ABACO.\newline \url{https://www.noticiasrcn.com/colombia/bancos-de-alimentos-activan-corredor-humanitario-tras-sismo-de-7-4-en-colombia-1048906} & Manizales como nodo de acopio; prioridades; puntos intermedios & 10 y 11 ago 2026 & 20 nov 2026 & no\\
Noticias RCN, ABACO supera 2.000 toneladas.\newline \url{https://www.noticiasrcn.com/colombia/abaco-supera-las-2-000-toneladas-de-ayuda-humanitaria-recolectadas-para-damnificados-por-el-terremoto-1051306} & Volumen del corredor; retos logísticos; bancos afectados & 18 ago 2026 & 20 nov 2026 & no\\
Portafolio, 2.000 toneladas movilizadas por ABACO.\newline \url{https://www.portafolio.co/historias-portafolio/abaco-2-000-toneladas-de-ayuda-movilizadas-502977} & Más de 9.470 familias en 36 municipios & 21 sep 2026 & 20 nov 2026 & no\\
La República, entrevista con el director ejecutivo de ABACO.\newline \url{https://www.larepublica.co/empresas/entrevista-con-juan-carlos-buitrago-director-ejecutivo-de-abaco-4458863} & Cadena de frío; prioridad a enlatados & Sin fecha en la investigación & 20 nov 2026 & no\\
Noticias RCN, 933 toneladas en el estadio Nemesio Camacho.\newline \url{https://www.noticiasrcn.com/colombia/ayudas-por-sismo-en-colombia-933-toneladas-en-el-nemesio-camacho-1051261} & Acopio de El Campín; 157 t para Caldas & Sin fecha en la investigación & 20 nov 2026 & no\\
Infobae, la Cruz Roja aclara qué donar.\newline \url{https://www.infobae.com/colombia/2026/08/14/cruz-roja-aclaro-que-donaciones-necesitan-las-victimas-del-terremoto-y-que-no-kits-de-habitat/} & Qué no donar; 280 t despachadas & 14 ago 2026 & 20 nov 2026 & no\\
Alcaldía de Bogotá.\newline \url{https://bogota.gov.co/mi-ciudad/gobierno/ayudas-de-bogota-ya-estan-en-quibdo-y-manizales-atencion-por-terremoto} & Tractomula con ayuda para Manizales & 15 ago 2026 & 20 nov 2026 & no\\
Publimetro, kits de alimentos del sector privado.\newline \url{https://www.publimetro.co/noticias/2026/09/11/asi-avanza-la-ayuda-para-las-familias-afectadas-por-el-terremoto-destinan-3000-millones-de-pesos-y-entregan-miles-de-kits-de-alimentos/} & 8.000 kits de mercado; entregas en Manizales, Supía y Riosucio & 11 sep 2026 & 20 nov 2026 & no\\
DANE, SIPSA, boletín semanal.\newline \url{https://www.dane.gov.co/files/operaciones/SIPSA/bol-SIPSASemanal-15ago21agos-2026.pdf} & Abastecimiento de la plaza Centro Galerías & 15 a 21 ago 2026 & 20 nov 2026 & no\\
Noticias RCN, desperdicio de alimentos según el DNP.\newline \url{https://www.noticiasrcn.com/colombia/colombia-desperdicia-9-76-millones-de-toneladas-de-alimentos-al-ano-1067616} & 9,76 millones de toneladas al año, 34~\% de la oferta & 3 oct 2026 & 20 nov 2026 & no\\
\midrule
\multicolumn{5}{l}{\textbf{Vías}}\\
Publimetro, cierres viales reportados por la Policía.\newline \url{https://www.publimetro.co/noticias/2026/08/10/piensa-viajar-por-carretera-tras-el-terremoto-estos-son-los-cierres-viales-en-colombia-reportados-por-la-policia/} & Cierre de la vía Pereira–Chinchiná & 10 ago 2026 & 20 nov 2026 & no\\
Semana, vías que siguen cerradas.\newline \url{https://www.semana.com/nacion/articulo/actualizacion-de-las-vias-que-siguen-cerradas-tras-el-terremoto-en-colombia/202602/} & Cierres vigentes el 12 de agosto & 12 ago 2026 & 20 nov 2026 & no\\
Colombia.com, estado de las vías.\newline \url{https://www.colombia.com/turismo/noticias/estado-de-las-vias-en-colombia-tras-el-sismo-del-10-de-agosto-2026-595846} & Cierres hacia Chinchiná y en el Alto de Letras & 10 y 12 ago 2026 & 20 nov 2026 & no\\
Minuto30, vías cerradas por el sismo.\newline \url{https://www.minuto30.com/vias-cerradas-por-sismo-colombia/1715765/} & Cierre de la vía Manizales–Fresno & 10 y 12 ago 2026 & 20 nov 2026 & no\\
ANI, corredores concesionados afectados.\newline \url{https://www.ani.gov.co/w/cuatro-corredores-viales-concesionados-presentaron-afectaciones-y-activaron-protocolos-de-seguridad-por-eventos-sismol%C3%B3gicos} & Derrumbes en el corredor Armenia–Pereira–Manizales & 10 ago 2026 & 20 nov 2026 & no\\
El Colombiano, vías afectadas.\newline \url{https://www.elcolombiano.com/negocios/vias-afectadas-temblor-quindio-caldas-valle-risaralda-JF39788045} & Daños típicos en las vías, según el INVÍAS & 10 ago 2026 & 20 nov 2026 & no\\
\midrule
\multicolumn{5}{l}{\textbf{Antecedentes y territorio}}\\
La FM, cierre de la vía a Bogotá en 2017.\newline \url{https://www.lafm.com.co/colombia/via-manizales-bogota-cerrada-derrumbe} & Cierre en el sector Brisas & 21 mar 2017 & 20 nov 2026 & no\\
El País, deslizamientos de abril de 2017.\newline \url{https://www.elpais.com.co/colombia/aumentan-a-17-los-muertos-tras-derrumbes-en-manizales-y-declaran-la-calamidad-publica.html} & 17 muertos y 75 viviendas colapsadas en barrios de ladera & 19 abr 2017 & 20 nov 2026 & no\\
La FM, reubicación tras 2017.\newline \url{https://www.lafm.com.co/colombia/cesar-uruena-pulido-tal-vez-sea-necesario-reubicar-500-familias-manizales} & Unas 500 familias por reubicar y 11 barrios en riesgo & 20 abr 2017 & 20 nov 2026 & no\\
El Tiempo, cierre del puente del Irra.\newline \url{https://www.eltiempo.com/colombia/otras-ciudades/cierre-indefinido-en-la-via-entre-manizales-y-medellin-174266} & Cierre preventivo de la vía hacia Medellín & 23 ene 2018 & 20 nov 2026 & no\\
ANI, cierre de la vía Manizales–Chinchiná.\newline \url{https://www.ani.gov.co/w/se-mantiene-el-cierre-total-de-la-v%C3%ADa-manizales-chinchin%C3%A1-debido-a-las-fuertes-lluvias} & Deslizamiento de unos 4.000~m³ y desvío por La Siria & 22 mar 2021 & 20 nov 2026 & no\\
Infobae, cierre de marzo de 2023.\newline \url{https://www.infobae.com/colombia/2023/03/16/cierre-total-en-la-via-manizales-bogota-por-deslizamientos/} & Cierre total e indefinido hacia Mariquita & 16 mar 2023 & 20 nov 2026 & no\\
La FM, cierre de junio de 2023.\newline \url{https://www.lafm.com.co/colombia/manizales-bogota-esta-completamente-cerrada} & Cierre en el Alto de Letras y desvíos & 7 jun 2023 & 20 nov 2026 & no\\
Infobae, cierre de diciembre de 2024.\newline \url{https://www.infobae.com/colombia/2024/12/30/bloqueada-la-via-manizales-bogota-por-un-derrumbe-autoridades-pidieron-paciencia-a-los-viajeros/} & Derrumbe cerca de Padua & 30 dic 2024 & 20 nov 2026 & no\\
Semana, reconocimiento de la ONU a Manizales.\newline \url{https://www.semana.com/mejor-colombia/articulo/reconocimiento-de-la-onu-a-manizales-por-su-modelo-de-gestion-del-riesgo/202626/} & Nodo de Resiliencia; alerta multiamenaza; intervenciones & 19 feb 2026 & 20 nov 2026 & no\\
Universidad Javeriana Cali, clústeres de afectación.\newline \url{https://www.javerianacali.edu.co/sites/default/files/2026-08/clu%CC%81steres%20de%20afectacio%CC%81n%20del%20sismo%20del%2010%20de%20agosto%20de%202026.pdf} & Manizales como clúster urbano de ladera & Sin fecha en la investigación & 20 nov 2026 & no\\
Alcaldía de Manizales, consolidado de juntas de acción comunal.\newline \url{https://manizales.gov.co/wp-content/uploads/2025/08/CONSOLIDADO-JUNTAS-ACCION-COMUNAL-Y-ADMINISTRADORAS-LOCALES-JAC-JAL.pdf} & 204 juntas, 12 comunas, 7 corregimientos, 131 barrios y 73 veredas & Sin fecha en la investigación & 20 nov 2026 & no\\
Telencuestas, proyección del DANE para 2025.\newline \url{https://telencuestas.com/censos-de-poblacion/colombia/2025/caldas/manizales} & Población de Manizales (fuente secundaria) & 2025 & 20 nov 2026 & no\\
\midrule
\multicolumn{5}{l}{\textbf{Plataformas, referentes y confianza}}\\
El Colombiano, Cuidar a Colombia (EAFIT).\newline \url{https://www.elcolombiano.com/inicio/cuidar-colombia-plataforma-eafit-estafas-terremoto-EK39884275} & Directorio verificado con 58 puntos de acopio en 32 ciudades & Sin fecha en la investigación & 20 nov 2026 & no\\
Noticias RCN, asistente digital para damnificados.\newline \url{https://www.noticiasrcn.com/economia/asistente-digital-gratuito-reune-informacion-oficial-para-damnificados-por-terremoto-1052956} & Bot de WhatsApp de JumpCube con AWS & Sin fecha en la investigación & 20 nov 2026 & no\\
Infobae, donaciones por Rappi.\newline \url{https://www.infobae.com/colombia/2026/08/11/asi-puede-donar-desde-20000-a-los-afectados-por-el-terremoto-usando-rappi/} & Botón de donaciones en Rappi & 11 ago 2026 & 20 nov 2026 & no\\
OPS, sistema SUMA/LSS.\newline \url{https://www.paho.org/es/temas/respuesta-emergencias-salud-desastres/sistema-lss} & Sistema regional de suministros humanitarios en la normativa colombiana & Sin fecha en la investigación & 20 nov 2026 & no\\
El País, inicio del Fondo Milagro.\newline \url{https://www.elpais.com.co/politica/de-la-espriella-anuncia-desde-pereira-el-inicio-del-fondo-milagro-para-reconstruir-viviendas-tras-el-terremoto-2215.html} & Objeto del Fondo Milagro & Sin fecha en la investigación & 20 nov 2026 & no\\
Informe de Robert Munro sobre Mission 4636.\newline \url{https://www.junglelightspeed.com/?p=745} & Mensajes de texto estructurados en Haití (2010) & Sin fecha en la investigación & 20 nov 2026 & no\\
Fast Company Brasil, Salva RS!\newline \url{https://fastcompanybrasil.com/impacto/aplicativo-gaucho-feito-em-menos-de-48-horas-ajuda-no-resgate-de-215-mil-pessoas/} & Web de pedidos de rescate en Rio Grande do Sul (2024) & Sin fecha en la investigación & 20 nov 2026 & no\\
El Financiero, Verificado19S.\newline \url{https://www.elfinanciero.com.mx/tech/asi-fue-como-verificado19s-puso-orden-a-la-ayuda-tras-el-sismo/} & Verificación de reportes en México (2017) & Sin fecha en la investigación & 20 nov 2026 & no\\
Semana, carrotanques de la UNGRD.\newline \url{https://www.semana.com/nacion/articulo/oxidados-golpeados-y-deteriorados-asi-estan-los-80-carrotanques-que-compro-la-ungrd-para-entregar-agua-en-la-guajira/202404/} & Antecedente de desconfianza en las ayudas & Sin fecha en la investigación & 20 nov 2026 & no\\
Noticias RCN, donaciones a la fundación de la primera dama.\newline \url{https://www.noticiasrcn.com/colombia/que-ha-pasado-con-las-donaciones-que-recibio-la-fundacion-de-la-primera-dama-tras-el-terremoto-1072443} & Preguntas públicas por las donaciones de 2026 & Sin fecha en la investigación & 20 nov 2026 & no\\
Herramientas de ruteo, registro y reportes (openrouteservice de HeiGIT, VROOM, Valhalla, LMMS, SCOPE, KoBo, Ushahidi). Síntesis en \ruta{docs/ESTADO_DEL_ARTE.md}; enlaces en \ruta{docs/investigacion/01_estado_del_arte_tecnico.md} & Qué existe y el hueco que ocupa DespachoX & Sin fecha en la investigación & 20 nov 2026 & no\\
\midrule
\multicolumn{5}{l}{\textbf{Documentos internos del equipo}}\\
\ruta{docs/PROBLEMA.md} & Descripción del problema, usuarios y ficha de la vertical 04 & 20 nov 2026 & 20 nov 2026 & no\\
\ruta{docs/investigacion/05_datos_abiertos_y_aliados.md} & Aliados posibles y qué pedirles; juntas de acción comunal & 20 nov 2026 & 20 nov 2026 & no\\
\ruta{docs/SIMULACION.md}, secciones 5 y 8, rama principal & Diseño y resultados de la simulación & 20 nov 2026 & 20 nov 2026 & no\\
\ruta{torre/CANAL.md}, entradas de las 4:57~p.\,m. del agente de simulación & Tres cifras para el pitch y su línea base; propuesta del peso 3; corrida con flota corta & 20 nov 2026 & 20 nov 2026 & no\\
\ruta{torre/CANAL.md}, entrada de las 3:55~p.\,m. & Primera prueba de punta a punta en la nube & 20 nov 2026 & 20 nov 2026 & no\\
\ruta{despacho/sim/resultados_s2/camperos_menos_2/resumen.md}, rama principal & Carga y número de semillas de la corrida con flota corta & 20 nov 2026 & 20 nov 2026 & no\\
\end{longtable}
}


\end{document}
